\section{Dynamical low-rank approximation}
\begin{figure*}[thbp]
    \centerline{\includegraphics[width=0.6\textwidth]{low_rank_manifold.png}}
    \caption{Illustration of the DLRA algorithm: the point $X^{(i)}$ lies on the manifold of rank-$r$ matrices $\mathcal{M}_r$, and $D^{(i)}$ is projected onto the tangent plane $\mathcal{T}_{X^{(i)}}(\mathcal{M}_r)$.}
    \label{fig:low_rank_manifold}
\end{figure*}
Consider the problem of finding a rank-$r$ solution $X$ via the update scheme $X^{(k+1)} = X^{(k)} + \alpha_k D^{(k)}$, where $D^{(k)}$ is a search direction and $\alpha_k > 0$ a step size. Starting from a rank-$r$ matrix $X^{(0)}$, there is in general no guarantee that $D^{(k)}$ preserves low rank, so that $X^{(k)}$ may leave the space $\mathcal{M}_r$ of rank-$r$ matrices after a single step.

One can replace $D^{(k)}$ by its projection onto the tangent space $\mathcal{T}_{X^{(k)}}(\mathcal{M}_r)$. This is the idea behind the Dynamical Low-Rank Approximation (DLRA) of Koch and Lubich \cite{kochDynamicalLowrankApproximation2007}. They proposed, in a continuous-time setting, evolving an approximation on $\mathcal{M}_r$ by projecting the derivative of the matrix onto the tangent space at the current approximation. Using this concept in our discrete setting, we replace $D^{(k)}$ by its tangent-space projection $\mathcal{P}_{X^{(k)}}[D^{(k)}]$, yielding the update
\begin{equation*}
    X^{(k+1)} = X^{(k)} 
    + \alpha_k \mathcal{P}_{X^{(k)}}(D^{(k)}).
\end{equation*}
This is illustrated in Figure~\ref{fig:low_rank_manifold}. Rather than updating $X^{(k)}$ directly, the DLRA algorithm operates on its SVD factorization $X^{(k)} = U \Sigma V^T$, updating the factors $U$, $\Sigma$, and $V$ at each iteration.

However, as pointed out by Koch and Lubich, when $\Sigma$ is ill-conditioned, which occurs when $r$ exceeds the effective rank of $X^{(k)}$, the projector $\mathcal{P}_{X^{(k)}}$ suffers from stiffness, leading to very small step sizes and issues with rounding errors.

To address this, Lubich and Oseledets proposed splitting the projector $\mathcal{P}_{X^{(k)}}$ into simpler sub-projections \cite{lubichProjectorsplittingIntegratorDynamical2013}. Crucially, each sub-projection can be evaluated without forming $\Sigma^{-1}$, and avoids the stiffness problem present in the original formulation. In recent years, additional advancements have been made to DLRA, offering improved stability, parallelism, and rank-adaptivity \cite{cerutiUnconventionalRobustIntegrator2020, cerutiRankadaptiveRobustIntegrator2021}.

Historically, DLRA was developed for time-continuous matrix differential equations. Schotthöfer et al. recently proposed a discrete rank-adaptive update scheme based on a robust projector-splitting integrator for training low-rank neural networks \cite{schotthoferLowrankLotteryTickets2022}. The discrete DLRA scheme used in this work follows their formulation.


\subsection{Application to the inverse problem}
In the source identification problem given in Eq.~\eqref{eq:problem}, the source $f \in L^2(\Omega)$ is often expected to be of low rank. If $x \in \mathbb{R}^{n^2}$ is some discretization of $f$, reordering $x$ into an $n \times n$ matrix $X = \operatorname{mat}(x)$, we expect that $r = \operatorname{rank}(X)$ is a small number, i.e., $r \ll n$. It is therefore natural to consider low-rank methods which can locate such solutions.

In the following sections, we consider the problem of minimizing a quadratic cost function
\begin{equation*}
    \Phi(X): \mathbb{R}^{n \times n} \rightarrow \mathbb R.
\end{equation*}
In our discrete case, the minimization problem given in Eq.~\eqref{eq:problem} translates into minimizing
\begin{equation*}
    \Phi(X) = \frac 1 2 \|Kx - y\|_{M_\partial}^2
    + \frac{\lambda^2}{2} \| W x \|_M^2.
\end{equation*}
Using the product rule, we get the gradient (with respect to $x$) $\nabla \Phi (X) = K^T M_\partial (K x - y) + \lambda^2 W^T M W x$, and the Hessian $H = K^T M_\partial K + \lambda^2 W^T M W$.


\subsection{Dynamical low-rank CG scheme}
Since $\Phi$ is quadratic with a symmetric positive definite (SPD) Hessian $H$ ($K^T M_\partial K$ is positive semi-definite and $\lambda^2 W^T M W$ is SPD), the conjugate gradient (CG) method is a natural minimization strategy \cite{saadIterativeMethodsSparse2003}.

CG generates a sequence of search directions $\{D^{(k)}\}$ that are $H$-conjugate, i.e.\ $\langle D^{(i)}, H D^{(j)}\rangle_F = 0$ for $i \neq j$. Starting from the initial search direction $D^{(0)} = -G^{(0)}$, each step computes an optimal step size
\begin{equation*}
    \alpha_k = \frac{\|G^{(k)}\|_F^2}{\langle D^{(k)},\, H D^{(k)}\rangle_F},
\end{equation*}
and updates the search direction by the Fletcher-Reeves rule \cite{fletcherFunctionMinimizationConjugate1964},
\begin{equation*}
    D^{(k+1)} = -G^{(k+1)} + \beta_k D^{(k)},
    \qquad
    \beta_k = \frac{\|G^{(k+1)}\|_F^2}{\|G^{(k)}\|_F^2}.
\end{equation*}

In our low-rank setting, we propose replacing the descent directions used in the method of Schotthöfer et al. with $H$-conjugate search directions $\{D^{(k)}\}$ obtained via the CG method. The motivation is that $H$-conjugate directions exploit curvature information of $H$, leading to faster convergence. The full scheme is given in Algorithm~\ref{alg:dlr-cg}.

We note that the ``Dynamical low-rank update'' part of Algorithm~\ref{alg:dlr-cg} is equivalent to Schotthöfer et al.'s method, with the search direction replaced. The $\operatorname{truncate}$-function is the rank-adaptive step of the algorithm, and takes in arguments \texttt{tol} and \texttt{max\_rank}.

\begin{algorithm}
\caption{Dynamical Low-Rank CG Scheme}
\label{alg:dlr-cg}
\begin{algorithmic}[1]

\Require Hessian $H$, initial matrix $X_0$ with SVD $X_0 = U_x \Sigma_x V_x^T$

\Ensure Low-rank solution $x = \operatorname{vec}(X)$

\State Compute initial gradient $G_0 = \nabla \Phi(X_0)$
\State Set initial search direction $D_0 = -G_0$

\For{$k = 0,1,2,\dots,\text{max\_iter}$}

    \State \textbf{Step size:}
    \State $\displaystyle \alpha_k = \frac{\|G_k\|_F^2}{\langle D_k, H D_k \rangle_F}$
    \Comment{$HD_k = \operatorname{mat}[H \operatorname{vec}(D_k)]$}

    \State \textbf{Dynamical low-rank update:}
    \State $W' = U_x \Sigma_x + \alpha_k D_k V_x$
    \State $L' = V_x \Sigma_x + \alpha_k D_k^T U_x$

    \State $\hat U = \operatorname{orth}([U_x, W'])$
    \State $\hat V = \operatorname{orth}([V_x, L'])$

    \State $\hat \Sigma = (\hat U^T U_x)\Sigma_x(V_x^T \hat V) + \alpha_k (\hat U^T D_k \hat V)$

    \State $(U_x, \Sigma_x, V_x) = \operatorname{truncate}(\hat U, \hat \Sigma, \hat V)$

    \State \textbf{Search direction update:}
    \State $G_{k+1} = G_k + \alpha_k H D_k$ 
    \Comment{or $G_{k+1} = \nabla \Phi(X_k)$}


    \State $\displaystyle \beta_k =
    \frac{\| G_{k+1} \|^2_F}
         {\| G_{k} \|^2_F}$
    \Comment{or $\displaystyle \beta_k = \frac{\langle G_{k+1}, H G_{k} \rangle_F}{\langle G_k, H G_k \rangle_F}$}
    \State $D_{k+1} = -G_{k+1} + \beta_k D_k$

\EndFor

\end{algorithmic}
\end{algorithm}


\subsection{Dynamical Low-Rank Preconditioned CG Scheme}
Preconditioning is a standard technique for accelerating CG methods, referred to as preconditioned CG (PCG). The system $H x = y$ is replaced by the equivalent system $P^{-1} H x = P^{-1} y$, where $P \approx H$ is the preconditioner matrix. The convergence of CG depends on $\kappa(H)$ (see Sec.~6.11.3 in Saad) \cite{saadIterativeMethodsSparse2003}, so $P$ is chosen such that $\kappa(P^{-1} H) \ll \kappa(H)$. For elliptic PDE-based forward operators $K$, the Hessian $H = K^T M_\partial K + \lambda^2 W^T M W$ typically has a large condition number, making preconditioning effective.

In the dynamical low-rank CG scheme, the preconditioned system $P^{-1} H x = P^{-1} y$ leads to the following adjustments in Algorithm \ref{alg:dlr-cg}:
\begin{equation*}
    \alpha_k = \frac
    {\langle G_k, Z_k \rangle_F}
    {\langle D_k, H D_k \rangle_F}, \qquad
    \beta_k = \frac
    {\langle G_{k+1}, Z_{k+1} \rangle_F}
    {\langle G_k, Z_k \rangle_F}, \qquad
    D_{k+1} = -Z_{k+1} + \beta_k D_k
\end{equation*}
These modifications are equivalent to standard preconditioned conjugate gradient, as originally described by Concus et al. \cite{concusGeneralizedConjugateGradient1976}. The resulting dynamical low-rank PCG scheme is summarized in Algorithm \ref{alg:dlr-pcg}.
\nomenclature[H]{$\Phi(X)$}{Cost function}
\nomenclature[H]{$H$}{Hessian of $\Phi(X)$}
\nomenclature[H]{$P$}{Preconditioner matrix}
\begin{algorithm}
\caption{Dynamical Low-Rank PCG Scheme}
\label{alg:dlr-pcg}
\begin{algorithmic}[1]

\Require Hessian $H$, initial matrix $X_0$ with SVD $X_0 = U_x \Sigma_x V_x^T$, preconditioner $P^{-1}$

\Ensure Low-rank solution $x = \operatorname{vec}(X)$

\State Compute initial gradient $G_0 = \nabla \Phi(X_0)$
\State Precondition $Z_0 = P^{-1} G_0$
\Comment $P^{-1} G_0 =
\operatorname{mat}[P^{-1} \operatorname{vec}(G_0)]$
\State Set initial search direction $D_0 = -Z_0$

\For{$k = 0,1,2,\dots,\text{max\_iter}$}

    \State \textbf{Step size:}
    \State $\displaystyle \alpha_k = \frac
    {\langle G_k, Z_k \rangle_F}{\langle D_k, H D_k \rangle_F}$
    \Comment{$HD_k = \operatorname{mat}[H \operatorname{vec}(D_k)]$}

    \State \textbf{Dynamical low-rank update:}
    \State $W' = U_x \Sigma_x + \alpha_k D_k V_x$
    \State $L' = V_x \Sigma_x + \alpha_k D_k^T U_x$
    \State $\hat U = \operatorname{orth}([U_x, W'])$
    \State $\hat V = \operatorname{orth}([V_x, L'])$
    \State $\hat \Sigma = (\hat U^T U_x)\Sigma_x(V_x^T \hat V) + \alpha_k (\hat U^T D_k \hat V)$
    \State $(U_x, \Sigma_x, V_x) = \operatorname{truncate}(\hat U, \hat \Sigma, \hat V)$

    \State \textbf{Search direction update:}
    \State $G_{k+1} = G_k + \alpha_k H D_k$
    \Comment{or $G_{k+1} = \nabla \Phi(X_k)$ }
    \State $Z_{k+1} = P^{-1} G_{k+1}$
    \State $\displaystyle \beta_k =
    \frac{\langle G_{k+1}, Z_{k+1} \rangle_F}
         {\langle G_k, Z_k \rangle_F}$
    \Comment{or $\displaystyle \beta_k = \frac{\langle G_{k+1}, H Z_{k} \rangle_F}{\langle G_k, H Z_k \rangle_F}$}
    \State $D_{k+1} = -Z_{k+1} + \beta_k D_k$

\EndFor

\end{algorithmic}
\end{algorithm}



\subsection{Choice of Preconditioner}
The effectiveness of PCG depends on the choice of $P$. A good preconditioner must approximate $H$ well, but must be cheap to apply at each iteration. We consider three common preconditioners: Jacobi scaling, incomplete Cholesky (IC), and (sparse) Cholesky. Since $H =  K^T M_\partial K + \lambda^2 W^T M W$ has a sparse plus low-rank structure, the Cholesky-based preconditioners are modified using the Woodbury formula to exploit this \cite{woodbury1950inverting}. We refer to Golub and Van Loan for a comprehensive treatment of preconditioning and iterative methods \cite{golubMatrixComputations4th2013}.

\paragraph{Jacobi scaling.} A Jacobi preconditioner $P = \operatorname{diag}(H)$ approximates $H$ by its diagonal. The approximation is cheap to compute, as $P^{-1}$ is the inverse of a diagonal matrix. It is most effective when $H$ is diagonally dominant.
 
\paragraph{Incomplete Cholesky.} For a symmetric positive definite (SPD) matrix $S$, the incomplete Cholesky factorization computes an approximation $S \approx  L L^T$, where $L$ is lower triangular. The Hessian $H = K^T M_\partial K + \lambda^2 W^T M W$ consists of a low-rank part plus a sparse part. We only apply the IC factorization to the sparse part $S = \lambda^2 W^T M W$, yielding $S \approx L L^T$. To account for the low-rank part, we let $K = U \Sigma V^T$ be the SVD of $K$ (which is assumed to be of rank $k$), and define $F = V \Sigma L_U$, where $L_U$ is the Cholesky factor of the $k \times k$ matrix $U^T M_\partial U$ (which is SPD since $M_\partial$ is SPD and $U$ has linearly independent columns). Then, the low-rank part can be written as $F F^T = V \Sigma L_U L_U^T \Sigma V^T = K^T M_\partial K$. By the Woodbury formula,
\begin{equation*}
    (S + FF^T)^{-1} = S^{-1} - 
    S^{-1} F (I_k + F^T S^{-1} F)^{-1} F^T S^{-1}.
\end{equation*}
Replacing $S^{-1}$ with the approximate inverse $L^{-T} L^{-1}$ yields the preconditioner
\begin{equation*}
    P^{-1} = L^{-T} L^{-1} -
    L^{-T} (L^{-1} F) C^{-1} (L^{-1} F)^T L^{-1},
\end{equation*}
where $C = I_k + (L^{-1} F)^T(L^{-1} F)$ is a $k \times k$ SPD matrix.

\paragraph{Sparse Cholesky.} The sparse Cholesky preconditioner follows the same Woodbury construction as the incomplete Cholesky preconditioner, but replaces the approximate factorization of $S$ with the exact sparse Cholesky factorization $S = LL^T$. Since the Woodbury formula then gives $P^{-1} = H^{-1}$, this is equivalent to solving $H x = y$ directly \textcolor{red}{(wrong)}, and CG converges in a single iteration. However, the setup is more expensive compared to incomplete Cholesky.
